% ============================================================
% Dịch vụ phân loại sản phẩm (Product Classification Service)
% Dùng trong: Chương 3 (Thiết kế) và Chương 4 (Cài đặt)
% ============================================================

% --- THIẾT KẾ (được \input từ analysis-design.tex) ---

\newcommand{\designClassify}{%

\section{Thiết kế dịch vụ phân loại sản phẩm}

Dịch vụ phân loại sản phẩm được xây dựng trên FastAPI 0.115, sử dụng PyTorch
2.5.1 và HuggingFace Transformers 4.46.3. Mô hình XLM-RoBERTa được load khi
khởi động thông qua FastAPI Lifespan Context Manager, tự động phát hiện GPU
và fallback về CPU nếu không có, đồng thời được đặt ở chế độ evaluation để
tối ưu hiệu năng suy luận.

Dịch vụ cung cấp hai endpoint phân loại chính: \texttt{POST /classify} cho
phân loại đơn lẻ với kết quả top-K predictions, và \texttt{POST /classify/batch}
cho xử lý hàng loạt tối đa 100 items với batch\_size 32. Ngoài ra, các endpoint
phụ trợ bao gồm \texttt{GET /categories} để lấy danh sách 32 danh mục,
\texttt{GET /model-info} để truy vấn metadata mô hình, và \texttt{GET /health}
để kiểm tra sức khỏe dịch vụ.

\begin{table}[htbp]
\centering
\caption{API endpoints dịch vụ phân loại}
\label{tab:classify-api}
\begin{tabular}{>{\raggedright\arraybackslash}p{0.35\linewidth} >{\raggedright\arraybackslash}p{0.55\linewidth}}
\toprule
\rowcolor{headerblue} \textcolor{white}{\textbf{Endpoint}} & \textcolor{white}{\textbf{Chức năng}} \\
\midrule
\rowcolor{rowgray} \texttt{POST /classify} & Phân loại 1 sản phẩm, trả top-K predictions \\
\texttt{POST /classify/batch} & Batch tối đa 100 items (batch\_size=32) \\
\rowcolor{rowgray} \texttt{GET /categories} & Danh sách 32 danh mục \\
\texttt{GET /model-info} & Metadata và metrics mô hình \\
\rowcolor{rowgray} \texttt{GET /health} & Health check \\
\bottomrule
\end{tabular}
\end{table}

Hệ thống phân loại sản phẩm vào 32 danh mục thương mại điện tử, bao gồm:
Electronics, Computers \& Networking, Mobile Phones \& Tablets, Fashion \& Clothing,
Shoes \& Footwear, Bags \& Luggage, Watches, Jewelry, Fashion Accessories,
Beauty \& Personal Care, Health \& Wellness, Sports \& Outdoors, Home Furniture,
Home Decor \& Lighting, Kitchen \& Dining, Bedding \& Bath, Large Appliances,
Small Appliances, Grocery \& Food, Baby \& Maternity, Toys \& Games, Books \& Media,
Stationery \& Office Supplies, Pet Supplies, Automotive \& Motorcycle,
Tools \& Hardware, Garden \& Outdoor Living, Musical Instruments,
Video Games \& Gaming, Software \& Digital Goods, Arts \& Crafts,
và Industrial \& Commercial.

}% end \designClassify

% --- CÀI ĐẶT (được \input từ implementation.tex) ---

\newcommand{\implClassify}{%

\section{Cài đặt dịch vụ phân loại sản phẩm}

\subsection{Pipeline suy luận}

\begin{enumerate}
\item Client gửi text mô tả sản phẩm qua POST request.
\item Validate input bằng Pydantic schema.
\item Tokenize input (max 512 tokens) qua AutoTokenizer.
\item Forward pass sinh logits cho 32 danh mục.
\item Softmax $\rightarrow$ top-K predictions với category ID, name, display name, score.
\item Trả response kèm inference time.
\end{enumerate}

Model được load khi khởi động qua FastAPI Lifespan Context Manager,
tự động detect GPU (fallback CPU), set evaluation mode.

\subsection{Hiệu năng}

\begin{table}[htbp]
\centering
\caption{Benchmark hiệu năng phân loại sản phẩm}
\label{tab:classify-benchmark}
\begin{tabular}{>{\raggedright\arraybackslash}p{0.35\linewidth} >{\centering\arraybackslash}p{0.2\linewidth} >{\centering\arraybackslash}p{0.2\linewidth}}
\toprule
\rowcolor{headerblue} \textcolor{white}{\textbf{Scenario}} & \textcolor{white}{\textbf{Latency}} & \textcolor{white}{\textbf{Throughput}} \\
\midrule
\rowcolor{rowgray} Single inference (CPU) & $\approx$ 45 ms & 22 req/s \\
Batch 10 items (CPU) & $\approx$ 180 ms & 55 items/s \\
\rowcolor{rowgray} Single inference (GPU) & $\approx$ 15 ms & 66 req/s \\
Batch 100 items (GPU) & $\approx$ 800 ms & 125 items/s \\
\bottomrule
\end{tabular}
\end{table}

}% end \implClassify
